\chapter{Henkin Completeness and Compactness}
\label{ch:fol-metatheory}

The completeness theorem for classical first-order logic is G\"odel's
\cite{Godel1930}.  This chapter develops the construction introduced by
Henkin \cite{Henkin1949}, which extracts a term structure from a maximal
consistent, witness-saturated set.

\section{Syntactic preparation}
\label{sec:fol-metatheory-01}

\begin{definition}[First-order consistency]
A set \(\Gamma\subseteq\Fm(\Lang)\) is \emph{consistent} iff
\(\Gamma\notderives_{\ND_{\FOL}}\bottom\).  It is \emph{maximal consistent}
iff it is consistent and \(\Gamma\cup\set{\varphi}\) is inconsistent for
every \(\varphi\notin\Gamma\).
\end{definition}

The closure properties in \cref{lem:prop-mcs} also hold for a maximal
consistent \(\Delta\subseteq\Fm(\Lang)\), but their first-order proof does not
silently graft derivations across eigenvariable rules.  For deductive closure,
suppose \(\Delta\derives_{\ND_{\FOL}}\varphi\) and
\(\varphi\notin\Delta\).  Maximality supplies a derivation of \(\bottom\) from
\(\Delta\cup\set{\varphi}\).  Because \(\Delta\) is consistent, at least one
open occurrence of the added assumption remains in that derivation.  Apply
\(\neg I\) at its root, discharging every such occurrence, to obtain
\(\Delta\derives_{\ND_{\FOL}}\neg\varphi\).  Combining this derivation with
the given derivation of \(\varphi\) by \(\neg E\) derives \(\bottom\) from
\(\Delta\), contradicting consistency.  Clauses 2--5 of
\cref{lem:prop-mcs} then use this deductive-closure result and only the
displayed propositional rules, so their proofs transfer without an
eigenvariable-renaming step.

\begin{lemma}[Classical quantifier duality]
\label{lem:quantifier-duality}
The calculus \(\ND_{\FOL}\) derives
\[
 \neg\forall x\,\varphi\leftrightarrow\exists x\,\neg\varphi
 \quad\text{and}\quad
 \neg\exists x\,\varphi\leftrightarrow\forall x\,\neg\varphi.
\]
\end{lemma}
\begin{proof}
For \(\neg\exists x\,\varphi\to\forall x\neg\varphi\), take a variable
\(y\) absent from \(\varphi\) and free in no open assumption.  Assume \(\varphi[y/x]\); existential
introduction contradicts \(\neg\exists x\varphi\).  Negation introduction
gives \(\neg\varphi[y/x]\), and universal introduction gives
\(\forall y\,\neg\varphi[y/x]\).  Since \(y\) is fresh for \(\varphi\), the
\(\alpha\) rule gives \(\forall x\neg\varphi\).  The reverse
implication follows by existential elimination: a witness to \(\varphi\)
contradicts the corresponding instance of \(\forall x\neg\varphi\).

For \(\exists x\neg\varphi\to\neg\forall x\varphi\), eliminate the
existential and contradict the corresponding universal instance.  For the
remaining direction, assume \(\neg\forall x\varphi\) and, for reductio,
\(\neg\exists x\neg\varphi\).  The already proved implication gives
\(\forall x\neg\neg\varphi\).  For a variable \(y\) absent from
\(\varphi\) and free in no open assumption, double-negation
elimination gives \(\varphi[y/x]\); universal introduction gives
\(\forall y\,\varphi[y/x]\), and the \(\alpha\) rule gives
\(\forall x\varphi\), contradicting the first assumption.  Reductio yields
\(\exists x\neg\varphi\).  Every eigenvariable is chosen absent from the
surviving assumptions and conclusion.
\end{proof}

\begin{lemma}[Fresh-constant lemma]
\label{lem:fresh-constant}
Let \(c\) be a constant symbol occurring in neither \(\Gamma\) nor
\(\varphi\).  If
\(\Gamma\derives_{\ND_{\FOL}}\varphi[c/x]\), then
\(\Gamma\derives_{\ND_{\FOL}}\forall x\varphi\).
\end{lemma}
\begin{proof}
Fix a derivation witnessing the premise.  Because the derivation is finite,
only a finite set \(\Gamma_0\subseteq\Gamma\) occurs among its open
assumptions.  Choose a variable \(y\) absent from \(\Gamma_0\), from
\(\varphi\), and from the finite derivation.  Uniformly replace \(c\) by \(y\)
in that derivation.  Induction on the derivation shows that every rule instance
remains an instance: \(c\) has no binding behavior, it occurs in no open
assumption, and \(y\) occurred in no original rule instance, so every
eigenvariable condition is preserved.  Hence
\(\Gamma_0\derives_{\ND_{\FOL}}\varphi[y/x]\).  Since \(y\) is free in no
member of \(\Gamma_0\), \(\forall I\) gives
\(\Gamma_0\derives_{\ND_{\FOL}}\forall y\,\varphi[y/x]\).  Because
\(y\notin\operatorname{Var}(\varphi)\), this conclusion is alpha-equivalent
to \(\forall x\varphi\); the \(\alpha\) rule therefore gives the latter.  The
same finite derivation also witnesses
\(\Gamma\derives_{\ND_{\FOL}}\forall x\varphi\).
\end{proof}

\begin{definition}[Witness property]
A set \(\Delta\subseteq\Fm(\Lang)\) has the \emph{witness property} iff, for
every formula \(\exists x\varphi\in\Delta\), there is a constant symbol \(c\)
of \(\Lang\) such that \(\varphi[c/x]\in\Delta\).
\end{definition}

The witness property is necessary in the term-structure construction: membership
of an existential formula must supply a term naming a witness.

\begin{lemma}[Witness-extension lemma]
\label{lem:witness-extension}
Let \(\Sigma\) be consistent, let \(c\) occur in neither \(\Sigma\) nor
\(\exists x\varphi\), and put
\[
 H_{x,\varphi,c}\defeq
 \exists x\varphi\to\varphi[c/x].
\]
Then \(\Sigma\cup\set{H_{x,\varphi,c}}\) is consistent.
\end{lemma}
\begin{proof}
Suppose the extension is inconsistent.  Negation introduction gives
\(\Sigma\derives\neg H_{x,\varphi,c}\).  Classical propositional derivations
then give
\[
 \Sigma\derives\exists x\varphi
 \quad\text{and}\quad
 \Sigma\derives\neg\varphi[c/x],
\]
because \(\neg(A\to B)\derives A\land\neg B\).  The fresh-constant lemma
applied to the second derivation gives
\(\Sigma\derives\forall x\neg\varphi\).  Quantifier duality gives
\(\Sigma\derives\neg\exists x\varphi\), contradicting the first derivation.
Thus the extension is consistent.
\end{proof}

\begin{lemma}[Eliminating newly added constants]
\label{lem:new-constant-elimination}
Let \(\Lang'\) extend \(\Lang\) only by constant symbols, and let
\(\Gamma\cup\set{\varphi}\subseteq\Fm(\Lang)\).  If there is an
\(\ND_{\FOL}\)-derivation of \(\varphi\) from \(\Gamma\) in the expanded
signature \(\Lang'\), then there is such a derivation in \(\Lang\).
\end{lemma}
\begin{proof}
Fix an expanded-signature derivation.  It is finite, so only finitely many
constants from \(\Lang'\setminus\Lang\), say
\(d_1,\ldots,d_n\), occur in it.  Choose distinct variables
\(z_1,\ldots,z_n\) absent from every formula and rule annotation in the
derivation.  Simultaneously replace each \(d_i\) by \(z_i\) throughout the
derivation.

Induction on the derivation verifies that every inference remains an instance
of the same rule.  Replacement commutes with the propositional constructors
and with term formation.  Because the \(z_i\) occurred in no original binder,
replacement causes no capture, and for every quantifier instance it commutes
with capture-avoiding substitution.  The eigenvariables of
\(\forall I\) and of \(\exists E\) are original variables,
distinct from all \(z_i\); replacing constants therefore preserves their
absence conditions.  Replacement also preserves alpha-equivalence and every
identity-rule instance.  The end formula and open assumptions contain none of
the \(d_i\), so they are unchanged.  The transformed tree is consequently a
\(\Lang\)-derivation of \(\varphi\) from \(\Gamma\).
\end{proof}

\begin{lemma}[Transfinite Lindenbaum lemma]
\label{lem:transfinite-lindenbaum}
Let \(\Lang'\) be a signature, let \(\lambda\) be an ordinal, and let
\(\langle\varphi_\alpha:\alpha<\lambda\rangle\) be an enumeration of
\(\Fm(\Lang')\) in which every formula occurs at least once.  Then every
set \(\Sigma\subseteq\Fm(\Lang')\) that is consistent for
\(\ND_{\FOL}\) is contained in a maximal consistent
\(\Delta\subseteq\Fm(\Lang')\).  The construction is choice-free given the
enumeration; the Axiom of Choice enters only when such an enumeration must
first be produced.
\textit{Status: complete proof.}
\end{lemma}
\begin{proof}
By \cref{thm:transfinite-recursion} with domain \(\lambda+1\), define a
sequence \(\langle\Gamma_\alpha:\alpha\le\lambda\rangle\) by
\[
 \Gamma_0=\Sigma,\qquad
 \Gamma_{\alpha+1}=
 \begin{cases}
   \Gamma_\alpha\cup\set{\varphi_\alpha},&\text{if this set is
     consistent},\\
   \Gamma_\alpha,&\text{otherwise},
 \end{cases}
 \qquad
 \Gamma_\beta=\bigcup_{\alpha<\beta}\Gamma_\alpha
 \ \text{ for limit }\beta.
\]

\emph{Chain claim:} if \(\alpha\le\beta\le\lambda\), then
\(\Gamma_\alpha\subseteq\Gamma_\beta\).  Proceed by transfinite induction
on \(\beta\): the claim is trivial at \(\beta=\alpha\), both successor
clauses extend the preceding stage, and a limit stage is the union of all
earlier stages.

\emph{Consistency:} every \(\Gamma_\alpha\) is consistent, by transfinite
induction on \(\alpha\).  Stage zero is the hypothesis on \(\Sigma\), and
the successor clause preserves consistency by its case split.  At a limit
\(\beta\), suppose \(\Gamma_\beta\derives_{\ND_{\FOL}}\bottom\).  The
derivation is finite, so its open assumptions form a finite
\(F\subseteq\Gamma_\beta\).  Each \(\psi\in F\) lies in some
\(\Gamma_{\alpha_\psi}\) with \(\alpha_\psi<\beta\).  If
\(F\ne\varnothing\), put
\(\alpha^{*}=\max_{\psi\in F}\alpha_\psi<\beta\), which exists because a
finite nonempty set of ordinals has a maximum; if \(F=\varnothing\), put
\(\alpha^{*}=0<\beta\).  The chain claim locates
\(F\subseteq\Gamma_{\alpha^{*}}\), so \(\Gamma_{\alpha^{*}}\) is
inconsistent, contradicting the induction hypothesis at \(\alpha^{*}\).

\emph{Maximality:} put \(\Delta=\Gamma_\lambda\); it is consistent, and it
contains \(\Sigma\) by the chain claim.  Suppose \(\varphi\notin\Delta\)
and pick \(\alpha<\lambda\) with \(\varphi_\alpha=\varphi\).  Rejection at
stage \(\alpha+1\) means that \(\Gamma_\alpha\cup\set{\varphi}\) derives
\(\bottom\).  A derivation from a set is a derivation from any superset,
because derivability from a set requires only that every open assumption
belong to it; hence \(\Delta\cup\set{\varphi}\) derives \(\bottom\).
Therefore \(\Delta\) is maximal consistent.
\end{proof}

\begin{lemma}[Henkin--Lindenbaum lemma]
\label{lem:henkin-lindenbaum}
Let \(\Lang\) be a signature and assume the Axiom of Choice in the
ambient set theory.  Every consistent
\(\Sigma\subseteq\Fm(\Lang)\) is contained in a maximal consistent set
\(\Delta\) in an expanded signature \(\Lang^H\), and \(\Delta\) has the
witness property.
\textit{Status: complete proof.}
\end{lemma}
\begin{proof}
 Let \(\kappa=\max(\abs{\Lang},\aleph_0)\).  Add new constant symbols
 \(\langle d_\xi:\xi<\kappa\rangle\).  Regard \(\Sigma\) as a set of formulas
 in the expanded signature.
It remains consistent there: otherwise
\cref{lem:new-constant-elimination} would transform an expanded-signature
derivation of \(\bottom\) into an original-signature derivation from
\(\Sigma\).  By \cref{cor:syntax-counting}, the expanded
language has at most \(\kappa\) formulas.  The \(\kappa\) formulas
\(\exists x_0\,(x_0=d_\xi)\), for \(\xi<\kappa\), are pairwise distinct, so
it has exactly \(\kappa\) formulas and exactly \(\kappa\) existential formulas.
Since a cardinal
is an ordinal by definition, fix a bijective enumeration
\[
  \langle E_\alpha:\alpha<\kappa\rangle,
  \qquad
  E_\alpha=\exists x_{n_\alpha}\varphi_\alpha
  \text{ with }n_\alpha\in\N,
\]
 of its existential formulas.  By the transfinite-recursion theorem
\cref{thm:transfinite-recursion}, at stage
\(\alpha<\kappa\) add a witness implication
\(H_{x_{n_\alpha},\varphi_\alpha,c_\alpha}\), where \(c_\alpha\) is the
least \(d_\xi\), in the displayed index order, that is absent from the current
theory and from \(E_\alpha\).  This least-available convention makes the
recursive clause functional.  Before stage \(\alpha\), fewer than
\(\kappa\) finite witness formulas have been added.  By
\cref{thm:infinite-cardinal-finite-coding}, their union contains fewer than
\(\kappa\) constants; \(E_\alpha\) contains only finitely many constants.
Therefore one of the \(\kappa\) new constants remains available.  The
witness-extension lemma
preserves consistency at successor stages.  At a limit stage take the union;
any derivation of \(\bottom\) would use finitely many assumptions and would
already occur at an earlier stage.  Let \(\Sigma^H\) be the resulting
consistent set.

The set \(\Fm(\Lang^H)\) has cardinality exactly \(\kappa\), as
established above from \cref{cor:syntax-counting} and the
formulas \(\exists x_0\,(x_0=d_\xi)\).  By
\cref{cor:cardinal-assignment}, under the Axiom of Choice already assumed,
fix a bijective enumeration
\(\langle\varphi_\alpha:\alpha<\kappa\rangle\) of \(\Fm(\Lang^H)\).
Apply \cref{lem:transfinite-lindenbaum} to \(\Sigma^H\) under this
enumeration to obtain a maximal consistent
\(\Delta\supseteq\Sigma^H\supseteq\Sigma\).  Every witness implication remains
in \(\Delta\).  If \(\exists x\varphi\in\Delta\), implication elimination
and deductive closure put its designated instance \(\varphi[c/x]\) in
\(\Delta\).  Hence \(\Delta\) has the witness property.
\end{proof}

